\section{Relation to previous work}

Nice to Meet You synthesizes practical MLIR abstract transformers and combines components using abstract meet~\cite{Peng2026,NTMYArtifact}. Its paper describes candidate soundness checking with Z3 across widths one through 64, discarding candidates on timeout. Our work studies a complementary route for a restricted residual language: replay establishes source equivalence and a product argument makes nine rows decisive for every positive width. It does not reproduce that system's synthesis coverage or replace its general solver use. The meaning of source operations must be matched explicitly before comparing verdicts; this paper's total word contract is not an implicit verification of the upstream compiler lowering.

Known-zero/known-one representations are established in LLVM~\cite{LLVMKnownBits}. The analysis of tristate numbers gives soundness and precision results for arithmetic transformers, including addition and subtraction~\cite{Vishwanathan2022}. We do not claim these abstract domains or their standard bitwise formulas as new. Here they provide a concrete finite semantic quotient, a complete classification of the specified coordinatewise fragment, and a target check attached to a source-bound certificate.

Bit-width-independent SMT encodings study a broader route through integer arithmetic, uninterpreted functions, and quantified constraints~\cite{Niemetz2019}. Our finite product argument is simpler because its residual expressions do not communicate between bit positions. It offers no general decision procedure for parametric bit-vector arithmetic. Alive2 instead checks LLVM transformations through bounded translation validation~\cite{Lopes2021}; its compiler refinement setting and treatment of the source language differ from the total transfer-expression contract used here.

The finite-ground congruence method underlying K1 is inherited from Paper II, building on the initial-model organization of Paper I~\cite{ShcherbakovI,ShcherbakovII}. The comparable-pair principal-congruence formula for semilattices used in Section~\ref{sec:sparse} is attributed to Martin-Maroto, Ricciardo, and de Polavieja~\cite{MartinMaroto2026}. This paper's contribution is the stated relation-matrix and sparse-representation contracts, together with the finite KnownBits classification and the proved composition of source reconstruction, local replay, and all-width target checking. The supplied evidence ties those claims to named implementations without identifying a tested program with a machine-checked proof.
